\documentclass[12pt]{article}
\usepackage[T1]{fontenc}
\usepackage[utf8]{inputenc}
\usepackage{lmodern}
\usepackage{textcomp}
\usepackage{enumitem}
\usepackage{geometry}
\geometry{margin=1in}
\begin{document}
\begin{center}
Answer Key - Economics - K-12 Level Assessment 6 \
\small (Answers are ordered exactly as in the question paper.)
\end{center}

\section*{Multiple Choice}
\begin{enumerate}[label=\arabic*.]
\item \textbf{Answer:} D
\\ \textit{Options:} A) an increase in the demand for fast-food workers.; B) a decrease in the supply of fast-food workers.; C) a shortage of fast-food workers.; D) a surplus of fast-food workers.
\\ \textit{Note:} A minimum wage set above the equilibrium wage creates a surplus of fast-food workers because it raises the cost of labor, leading to fewer jobs being available than there are workers willing to work at that wage.
\item \textbf{Answer:} C
\\ \textit{Options:} A) Labor; B) Capital; C) Human wants; D) Land
\\ \textit{Note:} Human wants are not considered a scarce economic resource because they are unlimited, while resources such as land, labor, and capital are finite and must be allocated efficiently to meet those wants.
\item \textbf{Answer:} D
\\ \textit{Options:} A) The demand for labor is independent of the demand for other inputs or resources.; B) The demand for labor is independent of the demand for the products produced by labor.; C) The demand for labor is independent of the availability of other inputs or resources.; D) The demand for labor is derived from the demand for the products produced by labor.
\\ \textit{Note:} The demand for labor is derived from the demand for the products produced by labor because employers hire workers to create goods and services only when there is a market demand for those products, reflecting the interdependent relationship between labor and product demand in the economy.
\item \textbf{Answer:} A
\\ \textit{Options:} A) appreciate.; B) depreciate.; C) remain unchanged.; D) change indeterminately.
\\ \textit{Note:} If interest rates rise more in country A than in country B, it attracts foreign investment, leading to increased demand for country A's currency, which causes it to appreciate.
\item \textbf{Answer:} A
\\ \textit{Options:} A) the number of firms increases; B) the price of output decreases; C) the labor supply curve shifts to the right; D) the labor supply curve shifts to the left
\\ \textit{Note:} The market demand curve for labor shifts to the right when the number of firms increases because more firms in the market create a higher overall demand for workers, leading to an increase in the quantity of labor demanded at each wage level.
\end{enumerate}

\section*{True / False}
\begin{enumerate}[label=\arabic*.]
\setcounter{enumi}{5}
\item \textbf{Answer:} True
\\ \textit{Options:} A) True; B) False
\\ \textit{Note:} The demand curve for a firm operating under perfect competition is perfectly horizontal because the firm is a price taker, meaning it can sell any quantity of its product at the market price but cannot influence that price due to the presence of many competitors offering identical products.
\item \textbf{Answer:} False
\\ \textit{Options:} A) True; B) False
\\ \textit{Note:} Normal goods can have either elastic or inelastic demand depending on factors such as the availability of substitutes and the proportion of income spent on the good, so it is incorrect to assert that they always have an elastic demand curve.
\item \textbf{Answer:} True
\\ \textit{Options:} A) True; B) False
\\ \textit{Note:} Buying Treasury securities from commercial banks increases the money supply in the economy, which is a fundamental action of expansionary monetary policy aimed at stimulating economic growth.
\item \textbf{Answer:} False
\\ \textit{Options:} A) True; B) False
\\ \textit{Note:} Increased consumer wealth and optimism typically lead to higher savings rates, which reduces the demand for loanable funds as consumers are less reliant on borrowing.
\item \textbf{Answer:} True
\\ \textit{Options:} A) True; B) False
\\ \textit{Note:} A monopoly typically results in more deadweight loss compared to a competitive market because it restricts output to raise prices above marginal cost, leading to decreased consumer surplus and overall market inefficiency.
\end{enumerate}

\section*{Long Form Response}
\begin{enumerate}[label=\arabic*.]
\setcounter{enumi}{10}
\item \textbf{Answer:} Increasing taxes and raising the discount rate are two strategies that can work together to combat inflation. When taxes are increased, consumers have less disposable income to spend, which can lead to reduced demand for goods and services. This decrease in demand can help lower prices and control inflation. Additionally, by raising the discount rate, the central bank makes borrowing more expensive, which can further reduce spending and investment in the economy. While these measures can effectively curb inflation, they may also pose challenges, such as slowing economic growth and leading to higher unemployment if consumers and businesses cut back too much on spending.
\\ \textit{Note:} correct because increasing taxes reduces consumer spending, while raising the discount rate discourages borrowing, both of which work together to decrease demand and help control inflation.
\item \textbf{Answer:} Matt's constant increase in total utility from consuming bratwurst would shape his demand curve to be horizontal, indicating that he is willing to purchase any quantity of bratwurst at a specific price. This means that as long as the price remains the same, Matt finds equal satisfaction in consuming more bratwurst, reflecting a consistent desire for the good. In economics, this behavior illustrates that consumers may have strong preferences for certain products, leading them to buy more without changing their willingness to pay. Consequently, it shows how consumer behavior can result in steady demand at particular price points, impacting market dynamics and pricing strategies.
\\ \textit{Note:} correct because it highlights that Matt's constant increase in total utility from bratwurst leads to a horizontal demand curve, indicating his willingness to buy any quantity at a constant price, which reflects the economic principle of strong consumer preference and consistent satisfaction from a good.
\end{enumerate}

\vfill
\noindent\textit{End of Answer Key}
\end{document}